\subsection{全集与映射}

让 \(\mathfrak{G}\) 是一个比集合论更强的理论，并令 \(\mathbf{E}\) 是...中的一项 \(\mathfrak{G}\) . 设 \(\Sigma\) 成为结构的一个种类，在 \(\mathfrak{G}\) 为简单起见，我们将在全文中假设 \(\Sigma\) 定义在单个（主）基集上，并且为简洁起见，我们将称“ \(\Sigma\) -集”为“赋予物种结构 \(\Sigma\) 的集合”。此外，我们假设 \(\sigma\) -态射已为物种定义 \(\Sigma\) ( \(\S 2\) ，第1号；如同在 \(\S 2\) 中，我们将用“态射”代替“ \(\sigma\) -态射”）。最后，物种 \(\Sigma\) 定义在基集上 \(x\) 并且具有 \(s\) 作为一般结构（ \(\S 1\) ，第4号），让我们假设一个项 \(\alpha \{x, s\}\) 在 \(\mathfrak{G}_{\Sigma}\) 中被定义，满足以下条件：

\((\mathbf{Q}\mathbf{M}_{\mathbf{I}})\) ，以及关系 \(\alpha \{x,s\} \subset \mathcal{F}(\mathbf{E};x)\) 在 \(\mathfrak{C}_{\Sigma}\) .

（QM \(_{II}\) ）如果在（一个理论 \(\mathfrak{G}'\) 该理论强于 \(\mathfrak{G}\) 中）F和F'是两个赋予结构的集合 \(\mathscr{S}\) , \(\mathscr{G}'\) 的物种 \(\Sigma\) ，且若 f 是 F 到 F' 的一个态射，则关系 \(\varphi \in \alpha \{F, \mathscr{G}\}\) 的极大性可知 \(f \circ \varphi \in \alpha \{F', \mathscr{G}'\}\) .

我们将这一关系表述为 \(\varphi\in\alpha\{x,s\}\) 通过说 \(\varphi\) 是 \(\alpha\) 从E到 \(x\) （赋予 \(s\) ).

一个 \(\Sigma\) -集合 \(\mathbf{F}_{\mathbf{E}}\) 和一个 \(\alpha\) -映射 \(\varphi_{\mathbf{E}}\) 的 \(\mathbf{E}\) 进入。 \(\mathbf{F}_{\mathbf{E}}\) 被称为万有的，如果满足以下条件：

(AU) 对每个 \(\alpha\) -映射 \(\varphi\) 的 \(\mathbf{E}\) 到 \(\Sigma\) -集合 \(\mathbf{F}\) 中存在唯一的态射 \(f\) 的 \(\mathbf{F}_{\mathbf{E}}\) 进入。 \(\mathbf{F}\) 使得 \(\varphi = f \circ \varphi_{\mathbf{E}}\) .

对 \((\mathbf{F}_{\mathbf{E}}, \varphi_{\mathbf{E}})\) 于是也称其为 E 的泛映射问题（相对于 \(\Sigma\) , \(\sigma\) ，以及 \(\alpha\) ).

让 \((\mathbf{F}_{\mathbf{E}}^{\prime},\varphi_{\mathbf{E}}^{\prime})\) 并且 \((\mathbf{F}_{\mathbf{E}}^{\prime \prime},\varphi_{\mathbf{E}}^{\prime \prime})\) 是 E 的泛映射问题的两个解。条件 (AU) 表明，存在唯一的态射 \(f_{1}\) 的 \(\mathbf{F}_{\mathbf{E}}^{\prime}\) 进入。 \(\mathbf{F}_{\mathbf{E}}^{\prime \prime}\) 和唯一的态射 \(f_{2}\) 的 \(\mathbf{F}_{\mathbf{E}}^{\prime \prime}\) 进入。 \(\mathbf{F}_{\mathbf{E}}^{\prime}\) 使得 \(\varphi_{\mathbf{E}}^{\prime \prime} = f_{1}\circ \varphi_{\mathbf{E}}^{\prime}\) 并且 \(\varphi_{\mathbf{E}}^{\prime} = f_{2}\circ \varphi_{\mathbf{E}}^{\prime \prime}\) 。因此我们有 \(\varphi_{\mathbf{E}}^{\prime} = f_{2}\circ f_{1}\circ \varphi_{\mathbf{E}}^{\prime}\) 并且 \(\varphi_{\mathbf{E}}^{\prime \prime} = f_{1}\circ f_{2}\circ \varphi_{\mathbf{E}}^{\prime \prime}\) 。将(AU)应用于 \(\mathbf{F} = \mathbf{F}_{\mathbf{E}}^{\prime}\) 并且 \(\varphi = \varphi_{\mathbf{E}}^{\prime}\) 的情形，我们得出 \(f_{2}\circ f_{1}\) 是 \(\mathbf{F}_{\mathbf{E}}^{\prime}\) 到自身的恒等映射。类似地， \(f_{1}\circ f_{2}\) 是 \(\mathbf{F}_{\mathbf{E}}^{\prime \prime}\) 到自身。因此（§2，第1号，准则CST8）， \(f_{1}\) 是……的一个同构 \(\mathbf{F}_{\mathbf{E}}^{\prime}\) 到上 \(\mathbf{F}_{\mathbf{E}}^{\prime \prime}\) ，以及 \(f_{2}\) 是它的逆同构。这个结果表述为：E的泛映射问题的解在同构意义下是唯一的。

为了验证一个对 \((\mathbf{F}_{\mathbf{E}}, \varphi_{\mathbf{E}})\) 是E的泛映射问题的一个解，通常方便的做法是验证以下两个条件：

（AU \(_{I}'\) ）对于每个Σ-集F和每个从E到F的α-映射φ，存在一个从F \(_{E}\) 到F的态射f，使得φ = f∘φ \(_{E}\) .

（AU \(_{II}'\) ）对于每个Σ-集F，两个从F \(_{E}\) 到F且在φ \(_{E}\) （E）上一致的态射相等。

因为如果这两个条件被满足，则由 \(f\) 所保证其存在的态射 \((\mathbf{AU}_{\mathbf{I}}^{\prime})\) 是唯一的，由 \((\mathbf{AU}_{\mathbf{II}}^{\prime})\) 。反之，显然(AU)蕴含 \((\mathbf{AU}_{\mathbf{I}}^{\prime})\) ；此外，如果 \(f\) 并且 \(f'\) 是 \(\mathbf{F}_{\mathbf{E}}\) 进入。 \(\mathbf{F}\) 的两个态射，它们在 \(\varphi_{\mathbf{E}}(\mathbf{E})\) ，我们有 \(f \circ \varphi_{\mathbf{E}} = f' \circ \varphi_{\mathbf{E}}\) ，由此 \(f = f'\) 上一致，则将(AU)应用于 \(\alpha\) -映射 \(f \circ \varphi_{\mathbf{E}}\) 。因此(AU)蕴含 \((\mathbf{AU}_{\mathbf{II}}^{\prime})\) .

